% 1-introduction.tex starts here
\section{Introduction}\label{cap:introducere}

\subsection{Motivation}\label{sec:motivatia-intro}
Food waste is a problem with a documented economic and social impact at both global and European level.
According to the \textit{UNEP Food Waste Index 2024} report~\cite{unepFoodWaste2024}, households generated over 631 million tonnes of wasted food in 2022 --- the equivalent of more than one billion meals thrown away every day, accounting for \(60\,\%\) of total food waste worldwide.
At the level of the European Union, Eurostat estimates that each inhabitant generated on average 132 kilograms of food waste in 2022, of which \(54\,\%\) came from households~\cite{eurostatFoodWaste2024}.
The European Parliament identifies among the main causes of waste at consumer level: impulse buying, the lack of adequate meal planning and poor management of the household food inventory~\cite{europarlFoodWaste2025}.
These behavioral causes are also confirmed by the academic literature: in a systematic review of 60 empirical studies on household food waste, Schanes et al.~\cite{schanes2018} identify shopping planning, checking the inventory before shopping and communication between household members as practices that play a decisive role in preventing over-purchasing.
The same review reports that simply using a shopping list reduced per-capita waste by approximately \(20\,\%\), and explicitly classifies applications that list the food inventory and plan meals among the technological prevention measures.
The authors stress, however, that the intention to reduce waste does not automatically translate into behavior (the attitude--behavior gap) and that some of the causes are structural, beyond individual control~\cite{schanes2018}; consequently, an application such as the one proposed in this thesis can help reduce redundant purchases without being a complete solution to the problem.

In parallel, modern \gls{web} applications raise significant technical challenges, documented both in the academic literature and in the annual surveys of the developer community.
The \textit{State of JavaScript 2024} survey identifies \emph{excessive complexity}, \emph{option overload} and \emph{breaking changes} as the main pain points reported by users of frontend frameworks~\cite{stateOfJS2024}.
The same source indicates a growing adoption of server-side rendering (\acrshort{ssr}) --- \(59\,\%\) of respondents use it in 2024 --- and a consolidation of \gls{typescript} as the standard language for large-scale applications: according to the \textit{JetBrains Developer Ecosystem 2024} survey, \gls{typescript} adoption grew from \(12\,\%\) in 2017 to \(35\,\%\) in 2024, and \(67\,\%\) of developers write predominantly \gls{typescript}~\cite{jetbrainsJsSurvey2024}.
These difficulties are amplified in full-stack applications, where poor integration of the layers leads to increased complexity and long-term maintenance costs.

This thesis does not treat these two directions as a relationship of necessity --- the food waste problem does not impose a particular architecture ---, but uses the first as a vehicle domain for investigating the second: a household application for pantry management, meal planning and shopping lists provides a representative full-stack context (integrated flows, derived data, relational persistence) in which the \gls{pan} architecture can be implemented and evaluated empirically.

\subsection{Purpose of the thesis}\label{sec:scop}
I set out to define and evaluate the \textbf{\gls{pan}} architecture, a full-stack \gls{web} model based on the \gls{typescript} ecosystem, which integrates the frontend, the backend and the persistence layer into a unified development process.
To investigate it empirically, I built \textit{\gls{kookio}}, a \gls{web} application that tackles the problem of household food waste.
Functionally, the application covers pantry management, daily meal planning and automatic derivation of the shopping list for household users.
Thus, \textit{\gls{kookio}} serves as a case study for evaluating the \gls{pan} architecture in a context representative of modern full-stack \gls{web} applications.

To this end, the thesis aims at:
\begin{itemize}
  \item implementing an integrated functional flow for the pantry, the daily planner and the derived shopping list, with persistence in a relational database;
  \item integrating the frontend and backend into a unified development process, with shared types and validation mechanisms;
  \item analyzing the characteristics and structural properties of the \gls{pan} architecture with respect to modularity (acyclic topology between libraries) and application maintainability (absence of a \textit{repository} layer);
  \item building a reusable \acrshort{cicd} pipeline, organized into \glspl{workflow} and \glspl{compositeaction}, aligned with the \gls{nx} targets;
  \item validating the implementation through unit tests with \gls{vitest} and \acrlong{e2e} (\acrshort{e2e}) with \gls{playwright}, together with an empirical evaluation of the persistence-layer latencies, front-end performance and test coverage.
\end{itemize}

\subsection{Contributions of the thesis}\label{sec:contributii}
\begin{enumerate}
  \item \textbf{An implemented \gls{pan} architectural model:} \gls{analog} (\gls{angular} 21) with server-side rendering (\acrshort{ssr}) on top of \gls{nitro}, \acrshort{trpc} for a typed API and \gls{prisma} to \gls{cockroachdb}, organized as an \gls{nx} \gls{monorepo} with four library groups and an acyclic topology enforced in \acrshort{ci} by the \gls{nx} \texttt{enforce-module-boundaries} rule.
  \item \textbf{The \textit{\gls{kookio}} application implemented as a \acrlong{mvp} (\acrshort{mvp}):} a working full-stack \gls{web} application for household users, which covers in an integrated way pantry management, daily meal planning, automatic derivation of the shopping list and the \textit{onboarding} flow (profile setup), with Firebase authentication, serving as a case study for the \gls{pan} architecture.
  \item \textbf{Patterns for sharing data contracts from a single source of truth:} the \gls{zod} schemas are generated automatically from \texttt{schema.prisma} by the \texttt{prisma-zod-generator} library~\cite{prisma_zod_generator}, and the \gls{typescript} types are derived by composition (\texttt{.pick()}, \texttt{.omit()}, \texttt{.extend()}, \texttt{z.infer<>}) from these schemas; the resulting contracts are reused in the user interface, in the server routes and in the tests, with no \textit{repository} layer between the \acrshort{trpc} procedures and the \gls{prisma} client.
  \item \textbf{A modular automated \gls{cicd} flow:} based on reusable \glspl{workflow} and \textit{\glspl{compositeaction}} aligned with the \gls{nx} targets, so that checks run only for the projects affected by a change (\texttt{nx affected}).
  \item \textbf{A testing methodology on two distinct levels:} unit testing with \gls{vitest} and \acrlong{e2e} (\acrshort{e2e}) with \gls{playwright}, together with concrete practices for the cold start of the \gls{vite} server (warm-up in \texttt{global-setup}) and for navigating to protected routes (\texttt{waitUntil: 'commit'}).
  \item \textbf{An empirical evaluation of the implemented technology stack:} measurements of latency at the persistence layer (\gls{cockroachdb} and \gls{redis}), of front-end performance (\gls{lighthouse} and image optimization) and of test coverage, summarized in \tabref{tab:rezultate}.
\end{enumerate}

The application is implemented as a working \gls{mvp}: recommending recipes based on the pantry inventory is already operational through progressive filtering (\secref{subsec:filtru-progresiv}), while the identified extensions aim at refining this recommendation with the dietary restrictions collected during \textit{onboarding} and at a mechanism for tracking staple consumables.
These extensions can strengthen the application's contribution to reducing food waste by cutting down redundant purchases.
The implementation details are presented in \secref{cap:implementare}, the empirical evaluation in \secref{ch:evaluare}, and the concrete directions for future development in \secref{sec:directii-viitoare}.
% 1-introduction.tex ends here
